\documentclass[10pt,a4paper]{amsart}
\usepackage{amsmath,amssymb,amsthm,graphicx}
\usepackage{lscape,longtable}

\def\be{\begin{equation}} \def\ee{\end{equation}}

\def\mrm{\mathrm}

%\def\ri{\mathrm{i}}
\def\ri{\mathrm{i}}

\def\bj{\mathbf{j}}

\def\fg{\mathfrak{g}}

\def\Zh{\widehat{Z}}

\def\sm{\mathsf{m}}

\def\BR{\mathbb{R}}
\def\BN{\mathbb{N}}
\def\BC{\mathbb{C}}
\def\BQ{\mathbb{Q}}
\def\BZ{\mathbb{Z}}
\def\BE{\mathbb{E}}
\def\BH{\mathbb{H}}

\def\CA{\mathcal{A}}
\def\CB{\mathcal{B}}
\def\CD{\mathcal{D}}
\def\CF{\mathcal{F}}
\def\CE{\mathcal{E}}
\def\CG{\mathcal{G}}
\def\CH{\mathcal{H}}
\def\CL{\mathcal{L}}
\def\CN{\mathcal{N}}
\def\CO{\mathcal{O}}
\def\CS{\mathcal{S}}
\def\wv{\vec{\rho}}

\def\ss{ \tilde{\mathcal{S}} }

\def\m{m_\ast}
\def\notowns{\not{\hspace{-0.7mm}\owns}}

\newcommand{\CC}{\mathbb{C}}
\newcommand{\RR}{\mathbb{R}}
\newcommand{\QQ}{\mathbb{Q}}
\newcommand{\FF}{\mathbb{F}}
\newcommand{\HH}{\mathbb{H}}
\newcommand{\ZZ}{\mathbb{Z}}
\newcommand{\NN}{\mathbb{N}}
\newcommand{\PP}{\mathbb{P}}
\newcommand{\MM}{\mathbb{M}}
\newcommand{\TT}{\mathbb{T}}
\newcommand{\Zhlong}[3]{\ensuremath{\widehat{Z}^{#1}_{#2}\left(#3;\tau\right)}}
\newcommand{\Zhshort}[3]{\ensuremath{\widehat{Z}^{#1}_{#2}\left(#3\right)}}
\newcommand{\Zchlong}[3]{\ensuremath{\widecheck{Z}^{#1}_{#2}\left(#3;\tau,\bar{\tau}\right)}}
\newcommand{\Zchshort}[3]{\ensuremath{\widecheck{Z}^{#1}_{#2}\left(#3\right)}}

\def\a{\alpha}
\def\ka{\vec \kappa}
\def\bon{\mathbf{n}}
%\def\bos{\vec{s}}
\def\bos{\varrho}
\def\sfs{\mathsf{s}}
\def\al{\alpha}
\def\bv{{\vec{\underline{b}}}}
\def\bal{\boldsymbol{\alpha}}
\def\bmu{\boldsymbol{\mu}}
\def\bb{\boldsymbol{\beta}}
\def\bol{\boldsymbol{\ell}}
 \def\bon{\boldsymbol{n}}
\def\brho{\boldsymbol{1}}
\def\bok{\boldsymbol{k}}
\def\boxx{\boldsymbol{x}}
\def\bosi{\bos}
\def\ex{{\rm e}}
\def\sgn{\mathrm{sgn}}
\def\tmod{\mathrm{mod}}
\def\id{\boldsymbol{1}}
\def\t{\tau}
\def\im{\operatorname{Im}}
\def\old{\text{old}}

\def\LII{{\rm II}}
\def\rr{ \mathsf{r}_\rho (\tilde{w}, \hat{w})}
\def\rp{ \mathsf{r}_\pi (\tilde{w}, \hat{w})}

%widecheck command definition
\DeclareFontFamily{U}{mathx}{\hyphenchar\font45}
\DeclareFontShape{U}{mathx}{m}{n}{
	<5> <6> <7> <8> <9> <10>
	<10.95> <12> <14.4> <17.28> <20.74> <24.88>
	mathx10
}{}
\DeclareSymbolFont{mathx}{U}{mathx}{m}{n}
\DeclareFontSubstitution{U}{mathx}{m}{n}
\DeclareMathAccent{\widecheck}{0}{mathx}{"71}


%Partial derivative definition
\makeatletter
\newcommand{\pder}[3][\@nil]{%
	\def\tmp{#1}
	\ifx\tmp\@nnil
	\ensuremath{\dfrac{\partial#2}{\partial#3}}
	\else
	\ensuremath{\dfrac{\partial^{#1}#2}{\partial#3^{#1}}}
	\fi}
\makeatother

%Regular derivative definition
\makeatletter
\newcommand{\der}[3][\@nil]{%
	\def\tmp{#1}
	\ifx\tmp\@nnil
	\ensuremath{\dfrac{d #2}{d #3}}
	\else
	\ensuremath{\dfrac{d ^{#1}#2}{d #3^{#1}}}
	\fi}
\makeatother

%Commands for repeating formulas (courtesy of Lucademicus, cf.\ https://tex.stackexchange.com/questions/459508/how-can-i-repeat-a-formula-without-copying-it)
\newcommand\labelAndRemember[2]
 {\expandafter\gdef\csname labeled:#1\endcsname{#2}%
 \label{#1}#2}
\newcommand\recallLabel[1]
 {\csname labeled:#1\endcsname\tag{\ref{#1}}}

%Comments
\newcommand{\mc}[1]{\textit{\color{red} MC: #1}}
\newcommand{\ic}[1]{\textit{\color{blue} IC: #1}}
\newcommand{\icf}[1]{\footnote{\it\color{blue} IC: #1}}
\newcommand{\gs}[1]{\textit{\color{green} GS: #1}}
\newcommand{\da}[1]{\textit{\color{orange} DP: #1}}

%Theorem environments
\newtheorem{thm}{Theorem}[section]
\newtheorem{prop}[thm]{Proposition}
\newtheorem{lem}[thm]{Lemma}
\newtheorem{conj}[thm]{Conjecture}
\newtheorem*{conj*}{General Conjecture}

\theoremstyle{definition}
\newtheorem{defn}[thm]{Definition}
\newtheorem{rmk}[thm]{Remark}

\numberwithin{equation}{section}

% the publisher class supplies the legacy font switch used by the author as {\cal X}; amsart does not
\providecommand{\cal}{\mathcal}
\usepackage{cleveref}
\crefname{def}{original equation}{original equations}
\Crefname{def}{Original equation}{Original equations}
\makeatletter
\newlabel{def:E01starless}{{internal}{source}{}{sourcepointer-1}{}}
\newlabel{def:E01starless@cref}{{[def][0][]internal}{[1][0][]source}{}{}{}}
\makeatother
\crefname{lem}{original lemma}{original lemmas}
\Crefname{lem}{Original lemma}{Original lemmas}
\makeatletter
\newlabel{lem:Estar_vs_starless}{{internal}{source}{}{sourcepointer-2}{}}
\newlabel{lem:Estar_vs_starless@cref}{{[lem][0][]internal}{[1][0][]source}{}{}{}}
\makeatother
\crefname{prop}{original proposition}{original propositions}
\Crefname{prop}{Original proposition}{Original propositions}
\makeatletter
\newlabel{prop:Estarless_eichler}{{internal}{source}{}{sourcepointer-3}{}}
\newlabel{prop:Estarless_eichler@cref}{{[prop][0][]internal}{[1][0][]source}{}{}{}}
\makeatother
\crefname{originalequation}{original equation}{original equations}
\Crefname{originalequation}{Original equation}{Original equations}
\makeatletter
\newlabel{Eichler_Enu_Com}{{internal}{source}{}{sourcepointer-4}{}}
\newlabel{Eichler_Enu_Com@cref}{{[originalequation][0][]internal}{[1][0][]source}{}{}{}}
\makeatother
\crefname{def}{original equation}{original equations}
\Crefname{def}{Original equation}{Original equations}
\makeatletter
\newlabel{def:Theta}{{internal}{source}{}{sourcepointer-5}{}}
\newlabel{def:Theta@cref}{{[def][0][]internal}{[1][0][]source}{}{}{}}
\makeatother
\crefname{prop}{original proposition}{original propositions}
\Crefname{prop}{Original proposition}{Original propositions}
\makeatletter
\newlabel{prop:E01asymp}{{internal}{source}{}{sourcepointer-6}{}}
\newlabel{prop:E01asymp@cref}{{[prop][0][]internal}{[1][0][]source}{}{}{}}
\makeatother
\crefname{s3}{original equation}{original equations}
\Crefname{s3}{Original equation}{Original equations}
\makeatletter
\newlabel{s3:hombl-2}{{internal}{source}{}{sourcepointer-7}{}}
\newlabel{s3:hombl-2@cref}{{[s3][0][]internal}{[1][0][]source}{}{}{}}
\makeatother
\crefname{eq}{original equation}{original equations}
\Crefname{eq}{Original equation}{Original equations}
\makeatletter
\newlabel{eq:abaction}{{internal}{source}{}{sourcepointer-8}{}}
\newlabel{eq:abaction@cref}{{[eq][0][]internal}{[1][0][]source}{}{}{}}
\makeatother
\crefname{lem}{original lemma}{original lemmas}
\Crefname{lem}{Original lemma}{Original lemmas}
\makeatletter
\newlabel{lem:shifting}{{internal}{source}{}{sourcepointer-9}{}}
\newlabel{lem:shifting@cref}{{[lem][0][]internal}{[1][0][]source}{}{}{}}
\makeatother
\crefname{originalequation}{original equation}{original equations}
\Crefname{originalequation}{Original equation}{Original equations}
\makeatletter
\newlabel{Section2-Fsplit}{{internal}{source}{}{sourcepointer-10}{}}
\newlabel{Section2-Fsplit@cref}{{[originalequation][0][]internal}{[1][0][]source}{}{}{}}
\makeatother
\title{Rewriting the companion function through non-holomorphic Eichler integrals}
\author{Cheng, Miranda C.N., Coman, Ioana, Passaro, Davide, Sgroi, Gabriele}
\date{doi:10.3842/SIGMA.2024.018; CC BY 4.0}
\begin{document}
\maketitle
\setcounter{section}{4}
\setcounter{thm}{1}
\noindent\textit{Editorial note.} The author mathematical text is preserved verbatim. The general proposition at1129-1143 is derived in its adjacent prose1123-1128, which explicitly combines four named earlier author results. This package retains that prose in its original position before the statement. The old package paired the general proposition to appendix2069-2108; that appendix begins with Brieskorn spheres and proves the specialization labelled lem:brieskorn\_thetas despite its subsection heading referring to lem:rewritingBBE. No mathematical argument has been supplied or rewritten.\par\medskip
\noindent\textit{Editorial note.} The retained preamble of this unit was generated mechanically for this delivery: it copies the author's own package, numbering, macro and theorem declaration block from the pinned source (source lines 9-156), with the publisher class and page setup lines removed; every package of that block is installed on the delivery host, and one line is appended defining the legacy font switch that the publisher class supplies and the author uses in the counted spans, which the delivery class does not provide.\par\medskip
\noindent\textit{Allowed standard prerequisites (editorial).}\begin{itemize}
\item The companion function of the rank two homological block. same article, the display at source line 1106 and the definitions at source lines 849 and 852. Fixes the object that is being rewritten and the block whose companion it is.
\item The non-holomorphic Eichler integral of a theta function and its properties. same article, the displays at source lines 893 and 898. Supplies the class of functions in which the rewritten expression is expressed and the identities used to collect the terms.
\item The asymptotic behaviour of the block and its companion. same article, the display at source line 770. Provides the asymptotic data with which the rewritten expression is compared in the last step.
\end{itemize}\medskip
\ifdefined\hypertarget\hypertarget{sourcepointer-1}{}\fi
\noindent\textit{Source reference.} Original equation internal: operations/codex-takeover-01a12558/corpus-repair-20261010/research-line/harness-ready-20261011/sources/sigma-2024-018-quantum-modular-widehat-z-g-invariants/pinned-original-source.tex, source lines 849--851. The counted statement or proof refers to that the author's definition of the starless block; it is the author's own earlier display, definition or statement in the same article and the whole argument of the counted span is retained verbatim.\par\medskip
\ifdefined\hypertarget\hypertarget{sourcepointer-2}{}\fi
\noindent\textit{Source reference.} Original lemma internal: operations/codex-takeover-01a12558/corpus-repair-20261010/research-line/harness-ready-20261011/sources/sigma-2024-018-quantum-modular-widehat-z-g-invariants/pinned-original-source.tex, source lines 852--856. The counted statement or proof refers to that the author's lemma comparing the two blocks; it is the author's own earlier display, definition or statement in the same article and the whole argument of the counted span is retained verbatim.\par\medskip
\ifdefined\hypertarget\hypertarget{sourcepointer-3}{}\fi
\noindent\textit{Source reference.} Original proposition internal: operations/codex-takeover-01a12558/corpus-repair-20261010/research-line/harness-ready-20261011/sources/sigma-2024-018-quantum-modular-widehat-z-g-invariants/pinned-original-source.tex, source lines 893--897. The counted statement or proof refers to that the author's proposition expressing the starless block through Eichler integrals; it is the author's own earlier display, definition or statement in the same article and the whole argument of the counted span is retained verbatim.\par\medskip
\ifdefined\hypertarget\hypertarget{sourcepointer-4}{}\fi
\noindent\textit{Source reference.} Original equation internal: operations/codex-takeover-01a12558/corpus-repair-20261010/research-line/harness-ready-20261011/sources/sigma-2024-018-quantum-modular-widehat-z-g-invariants/pinned-original-source.tex, source lines 898--902. The counted statement or proof refers to that the author's displayed Eichler integral of a theta function; it is the author's own earlier display, definition or statement in the same article and the whole argument of the counted span is retained verbatim.\par\medskip
\ifdefined\hypertarget\hypertarget{sourcepointer-5}{}\fi
\noindent\textit{Source reference.} Original equation internal: operations/codex-takeover-01a12558/corpus-repair-20261010/research-line/harness-ready-20261011/sources/sigma-2024-018-quantum-modular-widehat-z-g-invariants/pinned-original-source.tex, source lines 905--909. The counted statement or proof refers to that the author's definition of the theta functions used in the integrals; it is the author's own earlier display, definition or statement in the same article and the whole argument of the counted span is retained verbatim.\par\medskip
\ifdefined\hypertarget\hypertarget{sourcepointer-6}{}\fi
\noindent\textit{Source reference.} Original proposition internal: operations/codex-takeover-01a12558/corpus-repair-20261010/research-line/harness-ready-20261011/sources/sigma-2024-018-quantum-modular-widehat-z-g-invariants/pinned-original-source.tex, source lines 770--774. The counted statement or proof refers to that the author's proposition on the asymptotic behaviour of the block; it is the author's own earlier display, definition or statement in the same article and the whole argument of the counted span is retained verbatim.\par\medskip
\ifdefined\hypertarget\hypertarget{sourcepointer-7}{}\fi
\noindent\textit{Source reference.} Original equation internal: operations/codex-takeover-01a12558/corpus-repair-20261010/research-line/harness-ready-20261011/sources/sigma-2024-018-quantum-modular-widehat-z-g-invariants/pinned-original-source.tex, source lines 1106--1110. The counted statement or proof refers to that the author's display of the homological block whose companion is rewritten; it is the author's own earlier display, definition or statement in the same article and the whole argument of the counted span is retained verbatim.\par\medskip
\ifdefined\hypertarget\hypertarget{sourcepointer-8}{}\fi
\noindent\textit{Source reference.} Original equation internal: operations/codex-takeover-01a12558/corpus-repair-20261010/research-line/harness-ready-20261011/sources/sigma-2024-018-quantum-modular-widehat-z-g-invariants/pinned-original-source.tex, source lines 1734--1738. The counted statement or proof refers to that the author's display of the action of the Weyl group on the parameters; it is the author's own earlier display, definition or statement in the same article and the whole argument of the counted span is retained verbatim.\par\medskip
\ifdefined\hypertarget\hypertarget{sourcepointer-9}{}\fi
\noindent\textit{Source reference.} Original lemma internal: operations/codex-takeover-01a12558/corpus-repair-20261010/research-line/harness-ready-20261011/sources/sigma-2024-018-quantum-modular-widehat-z-g-invariants/pinned-original-source.tex, source lines 609--618. Navigation pointer to original author result retained in full source copy; this stub is not its proof.\par\medskip
\ifdefined\hypertarget\hypertarget{sourcepointer-10}{}\fi
\noindent\textit{Source reference.} Original equation internal: operations/codex-takeover-01a12558/corpus-repair-20261010/research-line/harness-ready-20261011/sources/sigma-2024-018-quantum-modular-widehat-z-g-invariants/pinned-original-source.tex, source lines 534--537. Navigation pointer to original author result retained in full source copy; this stub is not its proof.\par\medskip
In this subsection we will put the results obtained so far together and derive the form of the companion function for \smash{$\widehat Z^{{\rm SU}(3)}_{\bv}(M_3)$} for negative Seifert $M_3$ with three exceptional fibers, before we further specialize to the case of Brieskorn spheres.

Combining 1) (\ref{s3:hombl-2}),
the expression of \smash{$\widehat Z^{{\rm SU}(3)}_{\bv}(M_3)$} in terms of the generalized $A_2$ false theta function $F^{(\varrho)}$, 2) Lemma \ref{lem:shifting} and \eqref{Section2-Fsplit},
 the splitting of $F^{(\varrho)}$ into components, 3) Proposition~\ref{prop:E01asymp},
 the companion of the components, and 4) Proposition \ref{prop:Estarless_eichler} and Lemma \ref{lem:Estar_vs_starless}, the iterated non-holomorphic Eichler integral expressions for the companions, we finally obtain the following.

\begin{prop}
\label{lem:rewritingBBE}
For a negative Seifert manifold $M_3$ with three exceptional fibers, a companion function \smash{$\widecheck{Z}^{{\rm SU}(3)}_{\bv}(M_3)$} of the rank two homological blocks \smash{$\widehat Z^{{\rm SU}(3)}_{\bv}(M_3)$} is, up to potential one-dimensional pieces, given by the following non-holomorphic double Eichler integral
\begin{align*}% \label{Prop3-Zcheck}
\Zchlong{{\rm SU}(3)}{\bv}{M_3}
 ={}&z_{1\rm d}+ \frac{D}{m}\,C(q^{-1}) \sum\limits_{\hat{w}\in {\cal W}_\bv}(-1)^{\ell(\hat{w})} \sum_{\nu=0,1} \frac{\sqrt{3}}{4 \pi^{\nu}}	\left(\frac {3\Delta \vec{s}_{\hat w}}{m}\right)^{1-\nu}\\
 &\times\sum_{w\in W^+}\sum_{\delta\in\BZ/2} \bigl(\vartheta'_{w,\hat w,\delta},\vartheta^{1-\nu}_{w,\hat w,\delta}\bigr)^\ast(\tau),
\end{align*}
where the non-holomorphic double Eichler integral is of the theta functions
\begin{gather*}
 \vartheta'_{w,\hat w,\delta}(\tau)=
 \theta^1_{m, m\delta+\langle \vec\rho, w(\vec\sigma_{\hat w})\rangle}(\tau),\qquad
 \vartheta^{1-\nu}_{w,\hat w,\delta}(\tau) = \theta^{1-\nu}_{m,m\delta + \langle \Delta\vec\omega, w(\vec{\sigma}_{\hat w})\rangle}(3\tau).
\end{gather*}
\end{prop}

\end{document}
